\section{Komplexe Zahlen}

\begin{ejercicio}[Eine Rechenaufgabe mit komplexen Zahlen]
    Gegeben seien drei komplexe Zahlen $z_1 = 2 - i$, $z_2 = 5 + 2i$ und $z_3 = -3i$. Berechne:
    \begin{align*}
        \dfrac{\ol{z_1} + 2z_3 \cdot z_1 - iz_2}{2 (z_2 - \ol{z_3}) \cdot z_1}
    \end{align*}
\end{ejercicio}

\begin{ejercicio}[Rechnen mit komplexen Zahlen]~
\begin{enumerate}
    \item Berechnen Sie Realteil, Imaginärteil und Betrag der komplexen Zahl
    \begin{align*}
        z = 3 \cdot \ol{1+2i} \cdot (5 + i) + 5 \cdot \dfrac{28i - 3}{14 + (1 + 2i)^3}.
    \end{align*}

    \item Bestimmen Sie alle $z \in \bb{C}$, die folgende Gleichung lösen:
    \begin{align*}
        z \cdot \ol{z} + (3 - 4i) z + (3 + 4i) \ol{z} + 9 = 0.
    \end{align*}
    Beschreiben Sie die Gestalt der Lösungsmenge in der Gaußschen Zahlenebene.

    \item Berechnen Sie Realteil, Imaginärteil, Argument und Betrag von $z \in \bb{C}$:
    \begin{align*}
        z = 2^{-30} \cdot \left(i - \sqrt{3}\right)^{32}.
    \end{align*}

    \item Bestimmen Sie Betrag und Argument aller $v \in \bb{C}$, die im 2. Quadranten liegen (d.h. $\Re(z) \leq 0$, $\Im(z) \geq 0$) und folgende Gleichung lösen:
    \begin{align*}
        v^6 = -1 + i.
    \end{align*}

    \item Berechnen Sie
    \begin{align*}
        z = z_1 + z_2 + z_3 = 2 \cdot e^{-\nicefrac{5}{4} i \pi} + 0.2(1 - 3i)^3 + \dfrac{2 - i}{3 + 4i}.
    \end{align*}
    Geben Sie das Ergebnis sowohl in kartesischer als auch in exponenzieller Form an.
\end{enumerate}
\end{ejercicio}

\begin{ejercicio}[Beträge komplexer Zahlen]
    Beweisen oder widerlegen Sie: Für alle komplexen Zahlen gilt:
    \begin{align*}
        |z_1 + z_2|^2 + |z_1 - z_2|^2 = 2 \cdot \left(|z_1|^2 + |z_2|^2\right).
    \end{align*}
\end{ejercicio}